\begin{abstract}
La predicción de costos de seguros médicos es un problema de alto impacto
económico para aseguradoras y sistemas de salud. Este trabajo aplica
regresión supervisada sobre el Medical Cost Personal Dataset de Kaggle,
que contiene 1,338 registros con características demográficas y de salud.
Siguiendo la metodología CRISP-DM, se entrenaron y compararon seis modelos:
Regresión Lineal, Ridge, Lasso, Random Forest, XGBoost y Gradient Boosting.
Se incorporaron interacciones no lineales (BMI $\times$ tabaquismo, edad$^2$) como
ingeniería de features. El mejor modelo alcanzó un R$^2$ de [valor] y un MAE
de \$[valor], superando al baseline lineal en [X]\%. Los resultados
confirman el tabaquismo como el factor predictor más determinante del costo.

\textbf{Palabras clave:} regresión, XGBoost, Random Forest, seguros médicos,
predicción de costos, ingeniería de features, CRISP-DM.
\end{abstract}
